\documentclass[a4paper,10pt,notitlepage]{article}
\usepackage{ctex,geometry,graphicx,tikz,setspace,paralist,fancyhdr,caption}
\geometry{
	left=1.5cm,
	right=1.5cm,
	top=1.4cm,
	bottom=2cm,}
\newcommand{\rec}{
	\begin{tikzpicture}[remember picture,overlay]
		% 绘制边框
		\draw[line width=1.2pt] ([xshift=0.5cm,yshift=0.5cm] current page.south west) rectangle ([xshift=-0.5cm,yshift=-0.8cm] current page.north east);
	\end{tikzpicture}
	
}
\pagestyle{fancy}
\fancyhf{}
\fancyhead[C]{\rec}  % 在页眉中绘制图形
\fancyhead[L]{附录}
\fancyhead[R]{实验十 \quad 电压比较器}
\fancyfoot[C]{\thepage}
\begin{document}
	\large
	\onehalfspacing
	\begin{figure}[h]
		\raggedright
		\includegraphics{1.png}
	\end{figure}
	\centering
	{\Huge\textbf{模电实验报告}\par}
	\vspace{0.2cm}
	{\huge{实验内容：电压比较器}\par}
	\raggedright
	\vspace{0.3cm}
	\begin{centering}
		{\large 院系：电子与信息工程学院\hfill 学号：22309080\hfill 审批：\hspace{2cm} \par
			专业：通信工程\hfill 实验人：梁倍铭\hfill 日期：2023年12月14日 \par}
	\end{centering}
	\vspace{0.3cm}
	\section*{一、过零比较器}
	\begin{enumerate}
		\item 仿真结果：7.31V      实验结果：
		\item 仿真示意图如图9-1
		\begin{figure}[h]
			\centering
			\includegraphics[width=0.5\textwidth]{2.png}
			\caption*{图10-1}
		\end{figure}
		\item 改变$V_i$的幅值，$V_o$的幅值没有明显变化
	\end{enumerate}
	\section*{二、反相滞回比较器}
	\begin{enumerate}
		\item 仿真结果为：12.8V \qquad\qquad  实验结果为：
		\item 仿真结果为：12.8V \qquad\qquad  实验结果为：
		\item Vi 接500Hz 有效值lV 的正弦信号，观察并记录Vi 和Vo 波形。波形如图10-2
		\begin{figure}[h]
			\centering
			\includegraphics[width=0.5\textwidth]{3.png}
			\caption*{图10-2}
		\end{figure}
		\item 将电路中1MΩ 电位器调到200K，重复上述实验。\par
		由正到负 \qquad 仿真结果：12.8V   \qquad\qquad 实验结果：\par 
		由负到正 \qquad 仿真结果：12.8V   \qquad\qquad 实验结果：\par
		波形如图10-3
		\begin{figure}[h]
			\centering
			\includegraphics[width=0.5\textwidth]{4.png}
			\caption*{图10-3}
		\end{figure}
	\end{enumerate}
	\section*{三、同相滞回比较器}
	\begin{enumerate}
		\item 仿真结果为：14.8V \qquad\qquad  实验结果为：
		\item 仿真结果为：14.8V  \qquad\qquad 实验结果为：
		\item Vi 接500Hz 有效值lV 的正弦信号，观察并记录Vi 和Vo 波形。波形如图10-4
		\begin{figure}[h]
			\centering
			\includegraphics[width=0.5\textwidth]{5.png}
			\caption*{图10-4}
		\end{figure}
		\item 将电路中1MΩ 电位器调到200K，重复上述实验。\par 
		由正到负\qquad 仿真结果：13.8V  \qquad\qquad    实验结果：\par 
		由负到正\qquad 仿真结果：13.8V \qquad\qquad   实验结果：\par
		波形如图10-5
		\begin{figure}[h]
			\centering
			\includegraphics[width=0.4\textwidth]{6.png}
			\caption*{图10-5}
		\end{figure}
	\end{enumerate}
\end{document}